\documentclass[11pt]{article}
\usepackage[margin=24mm]{geometry}
\usepackage[T1]{fontenc}
\usepackage{lmodern,amsmath,amssymb,amsthm,mathtools,booktabs,microtype,hyperref}
\hypersetup{colorlinks=true,linkcolor=blue!50!black,urlcolor=blue!50!black}
\newtheorem{theorem}{Theorem}
\newtheorem{lemma}{Lemma}
\newcommand{\Z}{\mathbb Z}
\newcommand{\Q}{\mathbb Q}
\newcommand{\Fp}{\mathbb F_p}
\newcommand{\lcm}{\operatorname{lcm}}
\allowdisplaybreaks
\setlength{\emergencystretch}{2em}
\title{Adaptive CMF multiplication orders and a stronger arithmetic divisor\\
\large An addendum for the direction $(10,9,11,-11;2,1,3)$}
\author{Exact path experiments and a uniform arithmetic extension}
\date{10 October 2026}
\begin{document}
\maketitle

\begin{abstract}
We test literal changes of native CMF path, independently for each prime
and each prescribed initial row. Prime-dependent path selection strictly
improves finite divisibility certificates. Different rows can have
incompatible maximizing paths and a strictly better path-only divisor;
after adjoining the previous arithmetic normalization, separate row
choices give no further common-divisor gain in the tested families.
The residual prime pattern leads to a uniform extension of the previous
six-interval arithmetic lemma to all integer translates of its carry
chambers. Combined with the previously established analytic estimates,
the extension gives
$\mu(\zeta(2))\le62636575/13646619<4.589897$.
The path experiments and the uniform readout argument are distinguished
throughout: the new measure bound is not attributed to an unproved
row-dependent ordering rule.
\end{abstract}

\section{What was tested, and what improved}
The CMF is the native $\mathrm{pFq}(4,3,1)$ field, with
\[
w_0=(1,1,1,0;1,1,1),\quad v=(10,9,11,-11;2,1,3),\quad
R=\begin{pmatrix}0&2&-1\\2&0&0\end{pmatrix}.
\]
Write $W_N=W(w_0,w_0+Nv)$ and
$(p_N^{\rm raw},q_N^{\rm raw})^T=RW_Ne_0$.
Conservation identifies admissible paths with the same endpoint.
It does not permit unshifted matrices to be commuted.

For each cut vertex $w$, compute the complete native products
$A_w=W(w_0,w)$ and $B_w=W(w,w_0+Nv)$, and verify all nine entries
of $A_wB_w=W_N$ exactly. At this cut the certificate for initial row
$r$ and prime $p$ is
\[
b_{p,r}(w)=\min_{0\le i\le2}
 v_p\bigl((RA_w)_{r,i}(B_w)_{i,0}\bigr).
\]
It bounds the valuation of the endpoint scalar, since it bounds every
term in the three-term sum. Entire blocks are evaluated before taking
this minimum. The suffix always retains at least $2N$ negative-upper
moves; a cut immediately before the endpoint is excluded.

The two optimizations compared are
\[
 \max_w\min_r b_{p,r}(w),\qquad
 \min_r\max_w b_{p,r}(w).
\]
The first uses one path at each prime for both rows. The second also
allows a different path for each row. They can differ strictly, as the
example below shows. After adjoining the previous arithmetic certificate,
the resulting common-divisor bounds agree in these tests. This is a
finding about the tested certificates, not a general CMF theorem.

We used 89 structured cuts at each of
$N=20,40,80,160,320,640$, and 240 seeded random cuts at each of
$N=20,40,80,160,320,640,1280$. The random cuts also vary the progress
of the three lower parameters separately. All arithmetic in these
certificates is exact rational arithmetic. The 89-path search found
strict gains from prime-dependent selection. Beyond the previous
arithmetic normalization, the extra gain factor from row-dependent
selection was one at every tested depth. The random family gave the
same conclusion about that additional gain.

\subsection{A literal prime-dependent order witness}
At $N=80$, let path $A$ first increase each of $x_0,x_1,x_2$ by
$3N$, in that order; it then completes the remaining positive-upper
moves, the lower moves, and the negative-upper moves.
Path $B$ uses $4N$ instead of $3N$ for that initial block.
The cut is taken after the initial block. These are different
multiplication orders, since the remaining $x_0$ moves are performed
after some $x_1,x_2$ moves, with a different split in the two paths.
The suffix has all $11N$ negative-upper moves.
Their common-row valuation certificates are
\[
\begin{array}{c|rr}
 &p=2&p=59\\\hline
A&-19&-3\\
B&-25&-2
\end{array}
\]
Use $A$ at 2 and $B$ at 59. In the full 89-path search the sets of
maximizers at these two primes are disjoint. Thus this benefit is not
an artifact of choosing different names for tied optimal paths.
At $N=80$, prime adaptation improves the common factor over the best
single fixed path (after including the previous arithmetic bound)
by the factor $147649788$.

\subsection{A strict row-dependent order improvement}
At $N=40,p=491$, compare the cuts
\[
 w_A=w_0+N(3,3,3,-3;0,0,0),\qquad
 w_B=w_0+N(3,3,3,0;0,0,0).
\]
The first moves $3N$ negative-upper steps before completing the remaining
positive-upper and lower moves; the second postpones all negative-upper
moves until after those moves. Their certificates are
\[
\begin{array}{c|rr}
 &\text{numerator row}&\text{denominator row}\\\hline
A&-1&-2\\
B&-2&0
\end{array}
\]
In the entire 89-path family, the best common-path minimum is $-2$.
Choosing $A$ for the numerator and $B$ for the denominator gives $-1$.
Thus independent row choices improve the path-only common-divisor
certificate by a factor of $491$. The earlier arithmetic normalization
already certifies valuation $-1$ at this prime, so this gain does not
improve that normalization further.

For a second illustration, at $N=80,p=883$, a path with cut
\[
w=w_0+N(6,6,6,-9;2,1,3)
\]
has row certificates $(-2,-2)$, whereas the $3N$ initial-upper
path above has $(-3,0)$. The first better certifies the numerator
row, the second the denominator row. The sets of row maximizers in
the 89-path family are disjoint. Nevertheless the common minimum
remains $-2$: row adaptation adds no common factor in this example.
The numerator remains the limiting row in the aggregate experiments.

\subsection{Finite gains against the previous six-interval bound}
Each certificate is combined prime by prime with the earlier
normalization. The following are natural logarithms of upper bounds
on the \emph{reduced} denominator; the actual reduced denominator is
computed independently from the endpoint rational number.
\begin{center}
\begin{tabular}{rrrr}
\toprule
$N$&Previous bound&89-path bound&Actual $\log q_N$\\\midrule
20&1006.830013&969.294302&966.461089\\
40&2024.521715&1952.684660&1952.684660\\
80&3995.558516&3897.972570&3897.279423\\
160&7952.463445&7845.385158&7843.998864\\
320&15860.979328&15681.089752&15668.557979\\
640&31662.180409&31318.172451&31310.353216\\
\bottomrule
\end{tabular}
\end{center}
These finite gains include primes omitted by the cutoff $p^2>50N$;
that part contributes no exponential-rate improvement. It must not be
extrapolated as a constant saving per depth. The next section isolates
a uniform additional divisor among the eligible larger primes.

There is an important distinction: at $N=320,640$, every newly explained
large-prime factor is already exposed by the \emph{single} cut after all
positive-upper and lower moves, before the entire negative-upper block.
Different prime-dependent paths improve other finite factors, but the
new asymptotic divisor below does not require changing the path by prime.
It is therefore not evidence that adaptive orders are necessary for the
stronger irrationality measure.

\section{The uniform arithmetic extension}
This is an addendum to the supplied full manuscript
\emph{CMF conservation as an arithmetic tool}, saved as
\texttt{zeta2\_complete\_proof.tex} in the accompanying bundle.
We use its exact native readout, baseline normalization, eight-cell
signature table, and local-kernel lemma. Those statements include the
initial conditions and all native generator identities. The following
argument proves the extension needed here; it does not infer it from
the finite path tests.
\paragraph{A uniform periodic extension of the arithmetic divisor.}
The following statement extends the exact finite-field readout argument.
It is not a theorem asserting optimality, or necessity, of different
multiplication routes for the two initial rows.
Let
\[
 \mathcal T=\{(2,8,1),(3,9,2),(4,10,3),
                (5,8,3),(6,9,4),(12,9,8)\}.
\]
For $N\ge1$ and $r\in\{0,1,2\}$ define $H^*_{N,r}$ to be the product
of primes $p$ with $p^2>50N$ for which, with $t=\lfloor N/p\rfloor$,
at least one $(k,A,B)\in\mathcal T$ satisfies
\[
 \frac{13N+r}{13t+k}<p<\frac{AN}{At+B}.
 \tag{L1}
\]
Every prime occurs once. The inequalities imply $p<8N$, so the product
is finite. Its $t=0$ part is precisely the previous divisor $H_{N,r}$.
If there are no eligible primes, the product is one.

\begin{theorem}
With the same initial rows, trajectory, and factors $C_N,G_N$ as in
the preceding argument, for every integer $N\ge1$ and every
$r\in\{0,1,2\}$,
\[
 \boxed{\frac{C_Nd_{12N}d_{13N+r}}{G_NH^*_{N,r}}
             RW(N)e_r\in\mathbb Z^2.}
\]
\end{theorem}

\begin{proof}
The baseline integer normalization with $G_N$ has already been established.
It suffices to prove one additional factor $p$ for each prime in (L1).
Put $x=N/p=t+f$, where $0\le f<1$. The upper inequality in (L1) says
$f>B/A$; the lower says
\[
 f<\frac{k}{13}-\frac{r}{13p}.
\]
Thus $f$ lies in the same one of the six open intervals as in the
finite chamber table for the original shell argument.
The two intervals with an internal split are divided at $5/11$ and
$10/11$, exactly as before.

Write
\[
 a=\lfloor8x\rfloor=8t+a_0,\quad
 M=\lfloor11x\rfloor=11t+M_0,\quad
 b=8N-ap,\quad d=11N-Mp,\quad h=\phi(x)=\phi(f).
\]
The equality for $h$ follows directly from the periodic floor formula.
For $0\le j<p$, the carry indicators $e=\mathbf1_{j>d}$ and
$\delta_u=\mathbf1_{(j+uN)\bmod p<b}$ depend only on $f$ and $j/p$.
On their minimal-carry support they therefore have the same fixed
values as in the earlier table. The upper digits translate by
\[
 \lambda_u=\lfloor j/p+ux\rfloor=ut+\lambda_{u,0}
 \qquad (u=9,10,11).
\]
In particular they are still independent of the supporting residue $j$.
The two exceptional fractional split points cannot occur: if
$f=5/11$ or $10/11$, primality forces $p=11$ and
$N=11t+5$ or $11t+10$, contrary to $p^2>50N$.

Over $\mathbb F_p$ use the same polynomial quotient
\[
 \mathcal Q(X)=
 \frac{\prod_{u=9,10,11}\binom{X+uN}{b}}
 {(X^p-X)^{h-1}\prod_{j=0}^d(X-j)}.
\]
For $h=0$ the displayed formula means multiplication by $X^p-X$.
The carry multiplicities establish polynomial divisibility exactly as
before. Its decisive degree parameter is now
\[
 K=3a-M+h=13t+(3a_0-M_0+h)=13t+k.
\]
Hence (L1) gives
\[
 \deg(X^r\mathcal Q)=13N+r-(K-1)p-1<p-1,
 \qquad \sum_{j\in\mathbb F_p}j^r\mathcal Q(j)=0.
 \tag{L2}
\]

For clarity, the local kernel argument also remains valid with these
larger upper digits. Set
\[
 P(V)=\prod_{u=9,10,11}(a+1)^{\delta_u}
              \binom{\lambda_u+V}{a+\delta_u},\qquad
 U(V)=\frac{P(V)}{\prod_{l=0}^{M-e}(V-l)},\qquad
 \kappa=(-1)^{d-M+h-1}M!d!.
\]
The cutoff implies $p>50x$; the shell intervals imply $x>1/8$.
Thus $p$ is odd, $a+1<p$, and $M<p$, so all factorial denominators
and distinct kernel poles used here are legitimate over $\mathbb F_p$.
The exact separation of factors divisible by $p$ gives, as previously,
\[
 p^{1-h}R_{\rm ker}(j+pV)
       \equiv\kappa\mathcal Q(j)U(V)\pmod p,
\]
where $R_{\rm ker}(s)=F(s)(-1)^{11N}(11N)!/
\prod_{l=0}^{11N}(s-l)$ is the rational kernel, not the initial-row matrix.
Taking residues supplies the leading coefficients $c_l/p^h$.
Subtracting their partial fractions gives
\[
 p^{1-h}S_r(j+pV)
       \equiv\kappa j^r\mathcal Q(j)[U(V)]_+\pmod p,
 \tag{L3}
\]
with $[U]_+$ its polynomial part. This subtraction includes the
$e=1$ cases and the noncongruent contributions in the original
kernel. Residue classes outside the minimal-carry support give zero.
The previously established integer-valued normalization and degree
bound below $p^2$ justify (L3) also at its removable derivative nodes.

Choose the prime-dependent shift
\[
 \eta=p\left\lceil\frac{2N}{p}\right\rceil=Lp.
\]
If $p\le2N$, then $2N\le\eta<4N$; if $p>2N$, then
$\eta=p<8N$. Thus $2N\le\eta\le9N$ in every case, as required by
the exact readout. In the genuinely new periods $t\ge1$, one has
$p<N$ and the stronger bound $\eta<3N$.
All harmonic indices are below $20N<p^2$ (below $14N$ for $t\ge1$),
and the Newton degree plus one is at most $13N+r<p^2$.
Also $M+L\le13x+1<p$, so the high-digit harmonic reciprocals are units.
The integer factors at the relevant shifted Newton nodes have
absolute values below $24N+r<50N<p^2$.

The harmonic leading residue factors as a constant times the sum in
(L2), because the multiples of $p$ up to $j+(l+L)p$ are indexed by
$l+L$, independently of $j\in\{0,\ldots,p-1\}$.
For the polynomial part, (L3) and Lucas's congruence show that every
Newton coefficient with index $vp-1$ has the same factor
$\sum_jj^r\mathcal Q(j)$ and hence vanishes modulo $p$.
These are the only possibly nonintegral weights in the polynomial
functional. Both contributions therefore have valuation at least
$h-1$; the second row has valuation at least $h$.
Finally $v_p(d_{12N}d_{13N+r})=2$ and $v_p(G_N)=h$, so both
normalized entries acquire the asserted extra $p$.
Multiplying these prime divisibilities proves the statement for all $N$,
including the cases where the eligible product is empty.
\end{proof}

The arithmetic improvement in this theorem comes from the same exact
matrix readout and its prime-dependent polynomial analysis. Separate
finite experiments may expose these factors through specific
multiplication orders; their success alone is not the proof above.
\section{The improved rate}
\paragraph{Density of the lifted prime intervals.}
For
\[
 \mathcal B=\{(2,8,1),(3,9,2),(4,10,3),
                 (5,8,3),(6,9,4),(12,9,8)\},
\]
write $(k,a,b)$ for an element. The additional intervals, beyond the six
original intervals, are
\[
 \frac{13N+r}{13t+k}<p<\frac{aN}{at+b},\qquad t=1,2,\ldots .
\]
They are disjoint because their reciprocals lie in the disjoint translated
base intervals. If $H_{N,r}^{\mathrm{lift}}$ is the product of eligible primes
in these intervals, the prime number theorem gives
\[
 \lim_{N\to\infty}\frac{\log H_{N,r}^{\mathrm{lift}}}{N}
 =g:=\sum_{(k,a,b)\in\mathcal B}\sum_{t=1}^{\infty}
       \left(\frac{a}{at+b}-\frac{13}{13t+k}\right).
\]
To justify the infinite sum, first retain finitely many $t$. The column
shift changes each endpoint only by a bounded amount. The discarded
intervals for $t\ge T$ contain only primes $p<N/T$, so their total
logarithmic weight is at most $\vartheta(N/T)$; after division by $N$,
the limiting upper bound is $1/T$. The excluded primes
$p\le\sqrt{50N}$ contribute $o(N)$. This proves the stated limit without
requiring a prime number theorem uniform in $t$.
Equivalently,
\[
 g=\sum_{(k,a,b)\in\mathcal B}
       \bigl(\psi(1+k/13)-\psi(1+b/a)\bigr)
  =0.10868236921910157039\ldots.
\]

\paragraph{An exact rational enclosure.}
Each summand is
\[
 f_{k,a,b}(t)=\frac{ak-13b}{(at+b)(13t+k)}>0.
\]
Put $M=100000$ and $S=10^{24}$. Integer division gives
\[
 A:=\sum_{(k,a,b)\in\mathcal B}\sum_{t=1}^{M}
          \lfloor S f_{k,a,b}(t)\rfloor
   =108681309400407200714213.
\]
Since $f(t)=(\beta-\alpha)/((t+\alpha)(t+\beta))$ with
$\alpha=b/a<\beta=k/13$, comparison with decreasing reciprocal squares gives
\[
 \frac{ak-13b}{a(13M+13+k)}
 \le \sum_{t=M+1}^{\infty}f_{k,a,b}(t)
 \le \frac{ak-13b}{13aM}.
\]
Consequently
\[
 \frac AS+\sum_{\mathcal B}\frac{ak-13b}{a(13M+13+k)}
 \le g\le
 \frac{A+6M}{S}+\sum_{\mathcal B}\frac{ak-13b}{13aM}.
\]
The accompanying integer-arithmetic script verifies this enclosure:
\[
 0.1086823692136766<g<0.1086823692294671.
\]
Using the previously certified carry-density lower bound
\[
 \gamma\ge\frac{159010385641752609841}{13110552000000000000},
\]
set $c_0=25-233/120-\gamma_{\rm lower}$ and let $g_-$ denote the displayed
rational lower endpoint. The existing analytic enclosures imply
\[
 \sigma\le38.168732909647+c_0-g_-,\qquad
 \tau\ge24.467842131921-c_0+g_->0.
\]
Exact rational cross-multiplication then gives
\[
 1+\frac{38.168732909647+c_0-g_-}
          {24.467842131921-c_0+g_-}
   <4.589897.
\]
This last implication uses the uniform integer normalization furnished by
the translated shell argument and the analytic/nonvanishing inputs.
It does not follow merely from the finite divisibility checks below.

\paragraph{Why the column shift matters even after translation.}
At $N=3000$, $p=907$, the active lift has $t=3$, $k=4$ and
$(13t+k)p=43\cdot907=39001$.
Thus $13N+r<(13t+k)p$ holds only for $r=0$.
For the pair normalized using the original six-shell divisor, its two
$p$-valuations are $(1,3)$ in column zero, but $(0,2)$ in columns one and
two. The unshifted periodic guess therefore fails, whereas the shifted
inequality predicts precisely the observed distinction.
The same phenomenon occurs at $N=6000$, $p=1013$:
$(13\cdot5+12)1013=78001$.
\paragraph{The improved irrationality exponent.}
Let $H^*_{N,r}$ include all the translated shells proved in the
preceding arithmetic lemma. Their additional logarithmic density is
\[
 g=\sum_{(k,a,b)\in\mathcal B}\sum_{t\ge1}
 \left(\frac{a}{at+b}-\frac{13}{13t+k}\right),\qquad
 0.1086823692136766<g<0.1086823692294671,
\]
where
$\mathcal B=\{(2,8,1),(3,9,2),(4,10,3),(5,8,3),(6,9,4),(12,9,8)\}$.
The original six shells have density $233/120$. Together with the certified
carry-density lower bound, this gives
\[
 c=25-\gamma-\frac{233}{120}-g<10.821223.
\]
The integer pairs supplied by the arithmetic lemma consequently satisfy
\[
 |Q_N|\le\exp\bigl((48.989956+o(1))N\bigr),\qquad
 |Q_N\zeta(2)-P_N|
       \le\exp\bigl(-(13.646619-o(1))N\bigr).
\]
Here we used the analytic bounds $a<38.168733$ and
$b<-24.467842$. The previously established bounded-gap nonvanishing
property supplies the required nonzero rational comparison among three
consecutive indices. The standard irrationality-measure argument therefore
implies
\[
 \boxed{\mu(\zeta(2))\le
  1+\frac{48.989956}{13.646619}
  =\frac{62636575}{13646619}
  =4.589896955429033\ldots<4.589897.}
\]
All decimal inequalities here have exact rational certificates in the
accompanying code. The CMF approximants and their reduced denominators are
unchanged: the improvement is in the divisor that the argument certifies.

\section{Reproducibility and limits of the search}
The bundle includes the previous complete derivation, the present
addendum, all new search programs, exact checkpoints, the translated
kernel checks, and the rational density certificate. It also records
the eight-depth prime-pattern checks for all three output columns.
The pattern checks are numerical discovery evidence; the translated
finite-field argument above supplies the all-depth statement.

The requested new attachment
\texttt{zeta2\_cmf\_matrix\_proof (4).tex} was an empty file on disk.
The previous nonempty attachment was inspected as background. No
claim is made about the contents of the empty file.

The search changes arithmetic descriptions of the same endpoint,
so it cannot change the exact reduced approximants or their actual
delta values. No new delta search was performed. The previous values
were $\delta_{3000}=0.278673337629$,
$\delta_{4000}=0.279047437501$,
$\delta_{6000}=0.278946586656$, with last-three-depth slope estimate
$0.279152068816$. These numerical scores are separate from the
irrationality-measure conclusion proved using upper bounds.

No extra common-divisor benefit from independent row choice \emph{beyond
the previous arithmetic normalization} has been found in the tested
families. That remains a target for a richer lattice or residue
certificate. The precise next region is: keep the
same field, start, and direction; test $N=640,1280,2560$; retain at
least $2N$ negative-upper suffix moves; refine cut coordinates in
quarter-$N$ increments and allow one or two leading $p$-adic residue
digits. Examine primes outside the periodic divisor with $p^2>50N$.
At $N=6000$, remaining eligible residual primes include 1583 in all
three columns and 11743 in column 0. A new uniform rule is needed
before either can affect the stated measure bound.
\end{document}
